The regional channel-subset analysis above established how anatomical subsets of the montage behaved as self-contained detectors; the present analysis takes this reduction one step further and asks whether each \emph{individual} electrode of the full montage can independently support the complete Proposed-Med detector. For every electrode retained in the full-montage configuration, the complete Stage~A--Stage~C pipeline was rerun using that one electrode alone, with the same 30-second epochs, median centre, preprocessing and thresholding parameters as the full-montage reference, so each result is a genuine self-contained single-electrode detector rather than one electrode's score extracted from a multichannel run.

%As in the full-montage analysis (Section~\ref{sec:exp1_reference}) and the regional-subset analysis above, this comparison is restricted to the four anatomical regions, frontal, central, parietal and occipital, so 24 single-electrode detectors were evaluated independently on Internal and 18 on Cao2018 -- the same electrode set as Figure~\ref{fig:region_performance}, not the 32 electrodes of the complete montage nor the 6 (Internal)/4 (Cao2018) frontal electrodes examined in the earlier restricted single-channel analysis. One Internal electrode carries no conventional 10--20 label in the source hardware mapping and is reported under its raw EGI index, E83. As in every single-electrode selection, each result required no best-channel-per-session oracle and so corresponded directly to what a deployed single-electrode system would produce.
